\section{Discusión}

\subsection{Antipatrones en el dominio GenAI/RAG}

La literatura reciente coincide en que las aplicaciones que integran modelos generativos y mecanismos de recuperación acumulan formas específicas de deuda técnica. Shao et al. analizaron 100 aplicaciones LLM/RAG de código abierto, identificaron 18 patrones de defectos de integración y reportaron que el 77\,\% de las aplicaciones presentaba más de tres tipos de defectos \cite{shao_are_2024}. Sculley et al. habían mostrado antes que los sistemas de aprendizaje automático desarrollan deuda técnica oculta por interacciones entre datos, modelos, configuración y componentes de infraestructura \cite{sculley_hidden_2015}. Los estudios sobre antipatrones en sistemas de IA y aprendizaje automático amplían esa observación a nivel de diseño y de código \cite{bogner_characterizing_2021,gesi_code_2022,zhang_code_2022,oort_prevalence_2021,jebnoun_scent_2020,muralidhar_using_2021}.

El presente estudio se ubica en esa línea, pero no la replica. Shao et al. examinaron defectos funcionales de integración a partir de reportes y de una revisión orientada a la funcionalidad; este estudio examina antipatrones estructurales definidos operacionalmente y detectados mediante reglas AST sobre módulos Python. La comparación directa de porcentajes sería engañosa: el 77\,\% de Shao et al. se refiere a aplicaciones con más de tres tipos de defectos, mientras que el 91,3\,\% de este estudio se refiere a repositorios con al menos un hallazgo de cuatro antipatrones implementados. Son constructos, unidades y umbrales distintos. Lo que ambos resultados comparten es la dirección: el fenómeno no es marginal y afecta a una parte sustancial de los proyectos observados.

La arquitectura RAG añade una fuente adicional de complejidad: la recuperación introduce componentes, índices, embeddings y flujos de datos que deben coordinarse con la generación \cite{lewis_retrieval-augmented_2021,sharma_retrieval-augmented_2025,oche_systematic_2025}. Esa coordinación crea oportunidades para los antipatrones observados. AP1 concentra prompts en la lógica de control; AP3 fija proveedores y modelos; AP4 mezcla recuperación y generación en una sola función; AP5 silencia errores del modelo o del recuperador. Los cuatro son mecanismos plausibles de acoplamiento entre la capa de IA y la capa de aplicación. El estudio no demuestra que sean los mecanismos más frecuentes ni los más costosos, pero muestra que son detectables y que se concentran en módulos con peores métricas estructurales.

Los trabajos sobre deuda de seguridad y localización de defectos en aplicaciones con LLM refuerzan la pertinencia del dominio \cite{shen_security_2025,tan_lidl_2026}. También lo hacen las revisiones sobre el uso de LLM en ingeniería de software \cite{hou_large_2024}. Estas referencias no se utilizan aquí para validar la taxonomía, sino para situarla: el interés por la calidad estructural de las aplicaciones GenAI/RAG es reciente, creciente y todavía carente de instrumentos consolidados.

\subsection{Complejidad, mantenibilidad y acoplamiento}

La asociación observada entre antipatrones y mayor complejidad ciclomática, menor índice de mantenibilidad y mayor acoplamiento aproximado es coherente con la literatura clásica sobre olores y antipatrones. Garcia et al. propusieron los \emph{architectural bad smells} como señales de decisiones de diseño que pueden degradar la mantenibilidad \cite{garcia_identifying_2009}, y Fowler había formulado los \emph{code smells} como indicios de problemas de diseño que motivan refactorización \cite{fowler_1999}. Palomba et al. mostraron, en un estudio empírico a gran escala, que los olores de código se asocian con dificultades de mantenibilidad y con cambios más costosos \cite{palomba_diffuseness_2018}. Tufano et al. examinaron cuándo y por qué aparecen los olores, y señalaron que su introducción está ligada a decisiones de cambio y a presiones de entrega \cite{tufano_when_2015}.

El presente estudio no puede replicar ese tipo de evidencia longitudinal. Palomba et al. y Tufano et al. analizan la evolución del código y pueden establecer precedencia temporal entre la introducción del olor y el cambio; este estudio observa una sola versión por repositorio y, por tanto, no puede determinar si el antipatrón precedió al deterioro, si el deterioro favoreció la aparición del antipatrón o si ambos son consecuencia de una tercera característica del proyecto. La dirección de la asociación es consistente con la literatura, pero la inferencia causal no está justificada.

La distinción entre olores de código y antipatrones arquitectónicos también importa. Arcelli Fontana et al. mostraron que los olores arquitectónicos no son necesariamente independientes de los olores de código \cite{arcelli_fontana_are_2019}. En este estudio, los antipatrones se definen a nivel de módulo, pero se inspiran en preocupaciones arquitectónicas: separación de capas, portabilidad de proveedores, cohesión de funciones y propagación de errores. Esa mezcla de niveles es una fortaleza para la detección práctica y una debilidad para la interpretación teórica, porque un hallazgo puede reflejar tanto una decisión arquitectónica como un detalle de implementación.

El acoplamiento merece una advertencia específica. Chidamber y Kemerer definieron el CBO a nivel de clase, considerando las relaciones de dependencia entre clases \cite{chidamber_metrics_1994}. La métrica \texttt{cbo\_mod} utilizada aquí cuenta raíces de importación no estándar a nivel de archivo. Es una aproximación razonable para un análisis estático masivo, pero no captura herencia, composición, acoplamiento por parámetros ni dependencias dinámicas. El efecto mediano observado en el contraste de grupos y el efecto despreciable en el tercil superior deben interpretarse a la luz de esta definición restringida.

\subsection{Límites de la detección estática}

La literatura sobre detección de olores arquitectónicos ha documentado repetidamente que las herramientas heurísticas enfrentan un compromiso entre cobertura y precisión, y que sus resultados dependen de las reglas, los umbrales y el contexto del proyecto \cite{martini_identifying_2018,fontana_arcan_2017,mumtaz_systematic_2021}. Los catálogos de olores detectados por herramientas no siempre coinciden \cite{azadi_architectural_2019}, y la identificación de deuda arquitectónica requiere juicio experto además de automatización \cite{verdecchia_architectural_2018,de_toledo_identifying_2021}. Los trabajos sobre detección de olores en microservicios y en sistemas de aprendizaje automático confirman que la calibración del instrumento es un problema central \cite{pigazzini_towards_2020,gesi_code_2022}.

Los resultados de este estudio son consistentes con esa advertencia. La revisión asistida por IA de 60 hallazgos propuso 29 falsos positivos y 10 casos indeterminados. Aunque estas cifras no son una medida de precisión, sí indican que una parte importante de las emisiones podría corresponder a construcciones legítimas o ambiguas. AP3 es el caso más evidente: 15 de 17 casos revisados fueron propuestos como falsos positivos, lo que sugiere que la fijación literal de un modelo puede ser una práctica deliberada en adaptadores, ejemplos o configuraciones. AP1 y AP5 muestran una proporción más equilibrada, pero también con casos indeterminados.

La ausencia de una estimación de exhaustividad agrava el problema. Un detector conservador puede tener alta precisión y baja cobertura, o al contrario; sin un conjunto de referencia no es posible saberlo. El estudio no midió \emph{recall} y no puede afirmar qué proporción de los antipatrones presentes en el código fue detectada. Esta limitación es especialmente relevante para AP4, cuya regla exige más de sesenta líneas y dos tipos de llamadas en la misma función: es plausible que existan funciones monolíticas que no cumplen exactamente esos umbrales y que, por tanto, no se detectan.

\subsection{Adjudicación asistida por IA y validación humana}

El estudio distingue tres niveles que no deben confundirse. El primero es la detección automática: los cuatro detectores emiten hallazgos a partir de reglas AST. El segundo es la adjudicación asistida por IA: un modelo revisa 60 hallazgos muestreados y propone una clasificación con una justificación y un nivel de confianza. El tercero es la validación humana independiente: revisores humanos examinan el código y confirman o refutan cada hallazgo sin conocer la propuesta automática. Este estudio completó los dos primeros niveles y no completó el tercero.

La distinción no es retórica. La adjudicación asistida por IA puede ser útil para triar, priorizar y detectar patrones de error, pero no constituye un estándar de referencia. Un modelo puede compartir sesgos con el detector, aceptar una construcción plausible sin verificar el contexto completo o producir una justificación coherente para una clasificación incorrecta. Los 10 casos indeterminados son una señal de que la propia IA reconoce ambigüedad, y los 44 casos de confianza alta no equivalen a 44 aciertos. Sin un acuerdo entre revisores humanos y un cálculo de concordancia, no es posible afirmar que las etiquetas sean válidas.

La literatura metodológica sobre detección de olores insiste en la necesidad de validación empírica y de acuerdo entre evaluadores \cite{martini_identifying_2018,mumtaz_systematic_2021}. Este estudio no cumple ese estándar. Por ello, la discusión de los resultados debe mantenerse en el plano de las hipótesis y de las asociaciones, y no en el de las conclusiones validadas. La revisión asistida por IA se conserva como un antecedente útil para diseñar una validación humana posterior, no como un sustituto de ella.

\subsection{Implicaciones para la práctica}

A pesar de las limitaciones, los resultados tienen implicaciones prácticas acotadas. Primero, sugieren que los antipatrones definidos son lo bastante frecuentes como para justificar una revisión sistemática en proyectos GenAI/RAG. Segundo, identifican AP1 y AP5 como los fenómenos más extendidos en la muestra: prompts incrustados en la lógica de control y errores críticos silenciados. Tercero, muestran que la concentración por repositorio es alta, de modo que una auditoría puede comenzar por los proyectos con mayor densidad de hallazgos y por los módulos con más de un antipatrón.

La aplicación práctica debe ser prudente. El instrumento puede usarse como tamiz para orientar una revisión humana, no como puerta de calidad automática. Un hallazgo de AP1 puede ser un prompt legítimo que conviene externalizar; un hallazgo de AP3 puede ser una decisión deliberada de despliegue; un hallazgo de AP5 puede ser un manejo de errores aceptable en un contexto específico. La herramienta señala candidatos; la decisión requiere contexto, juicio y, idealmente, revisión por pares.

Desde la perspectiva de la arquitectura RAG, los resultados sugieren tres recomendaciones provisionales: separar la recuperación de la generación en funciones y módulos con responsabilidades claras, externalizar prompts y configuración de proveedores para reducir el acoplamiento, y revisar los manejadores de excepciones que silencian fallos del modelo o del recuperador. Estas recomendaciones son hipótesis de diseño derivadas de la taxonomía, no prescripciones validadas por un experimento controlado.

\subsection{Aportes y relación con el estado del arte}

El estudio aporta cuatro elementos. Primero, una taxonomía operacional de seis antipatrones específicos para aplicaciones GenAI/RAG, con reglas AST, atributos afectados y exclusiones. Segundo, un instrumento de análisis estático reproducible que se ejecutó sobre 23 repositorios y 15\,769 archivos. Tercero, una caracterización empírica de 445 hallazgos en 201 módulos, con distribución por antipatrón y por repositorio. Cuarto, un análisis estadístico que reporta tamaños de efecto, correlaciones, terciles de tamaño y sensibilidad por repositorio, junto con una revisión asistida por IA de 60 casos.

En relación con el estado del arte, el estudio complementa los trabajos orientados a defectos funcionales de integración \cite{shao_are_2024,tan_lidl_2026} con una mirada estructural y arquitectónica. También dialoga con la literatura sobre deuda técnica en sistemas de IA \cite{sculley_hidden_2015,bogner_characterizing_2021,gesi_code_2022} y con la tradición de detección de olores \cite{garcia_identifying_2009,fowler_1999,palomba_diffuseness_2018}. Su contribución no es reemplazar esos trabajos, sino explorar un ángulo poco cubierto: la detección automatizada de antipatrones arquitectónicos en el código de aplicaciones GenAI/RAG y su asociación con métricas tradicionales de calidad.

\subsection{Comparabilidad y validez externa}

La comparación con otros estudios debe hacerse con cautela por cuatro razones. Primera, la muestra es no probabilística y está compuesta por 23 repositorios, con 13 en la banda de más de 1000 estrellas; no representa al conjunto de aplicaciones GenAI/RAG. Segunda, el análisis se restringe a Python y a código open source; no cubre aplicaciones propietarias, otros lenguajes ni sistemas que integran modelos mediante servicios externos sin código disponible. Tercera, los constructos difieren: un defecto funcional de integración no es lo mismo que un antipatrón estructural, y una aplicación no es lo mismo que un módulo. Cuarta, la concentración de hallazgos en dos repositorios hace que los promedios globales dependan de esos proyectos.

Estas limitaciones no impiden extraer lecciones, pero impiden generalizar. El valor del estudio es el de un piloto metodológico: muestra que el instrumento funciona, que los fenómenos son observables y que las asociaciones merecen una investigación más rigurosa. La generalización requerirá muestras más amplias, estratificación explícita, modelos multinivel y validación humana independiente.
